// Break Search — scan for the first bright pixel using for + break
//
// Educational example demonstrating for loops, break, and fetch().
// For each row (iy), scans horizontally from the left edge to find the
// first pixel whose luminance exceeds the threshold. Outputs a
// visualization: pixels left of the found position are tinted red,
// and the found column is drawn as a bright green line.
//
// The brightness test reads the image, so it is per-pixel: a `break` under
// it would end the scan for EVERY pixel the first time the loop reached it
// (LANGUAGE.md §7.1). Each pixel records its own first hit in found_x
// instead, and `break` is kept for a test that is the same for every pixel:
// the scan stops at the right edge of the image.

f$threshold = 0.5;    // luminance threshold to search for
i$scan_width = 0;     // pixels to scan (0 = full width)

vec3 src = @image;

// Determine how far to scan
int width = ($scan_width > 0) ? $scan_width : int(iw);

// Scan this row from the left to find the first bright pixel
int found_x = -1;
for (int sx = 0; sx < width; sx = sx + 1) {
    // Uniform: past the last column there is nothing left to find
    if (sx >= int(iw)) {
        break;
    }
    vec3 s = fetch(@image, sx, iy);
    float lum = luma(s);
    // Per-pixel: a pixel that already found its hit keeps it
    if (found_x < 0 && lum > $threshold) {
        found_x = sx;
    }
}

// Visualize the result
vec3 result = src;
if (found_x >= 0) {
    if (ix < found_x) {
        // Tint pixels before the found position red
        result = vec3(src.r * 0.5 + 0.3, src.g * 0.3, src.b * 0.3);
    } else if (ix == found_x) {
        // Draw a bright green line at the found position
        result = vec3(0.0, 1.0, 0.0);
    }
}

@OUT = result;
